% Part: first-order-logic
% Chapter: completeness
% Section: identity

\documentclass[../../../include/open-logic-section]{subfiles}

\begin{document}

\olfileid{fol}{com}{ide}
\olsection{Identitas}

\begin{explain}
Konstruksi model term ing bagean sadurunge cukup kanggo netepake
kelengkapan logika orde kapisan tumrap himpunan~$\Gamma$ kang ora
ngemot~$\eq$. Model term kasebut nyukupi saben $!A \in \Gamma^*$ kang
ora ngemot~$\eq$ (lan mula kabeh $!A \in \Gamma$). Nanging konstruksi
iki ora lumaku yen $\eq$ ana. Sebabe, $\Gamma^*$ banjur bisa ngemot
!!a{sentence}~$\eq[t][t']$, nanging ing model term, nilai saben term
yaiku term kasebut dhewe. Mula, yen $t$ lan $t'$ iku term kang beda,
nilai kalorone ing model term---yaiku $t$ lan $t'$---uga beda, saengga
$\eq[t][t']$ luput. Nanging iki bisa didandani nganggo konstruksi kang
dikenal minangka ``pemfaktoran.''
\end{explain}

\begin{defn}
  Upamana $\Gamma^*$ minangka himpunan ukara kang konsisten lan
  !!{complete} ing~$\Lang L$. Kita nemtokake relasi $\approx$ ing
  himpunan term katutup saka~$\Lang L$ kanthi
  \[
  t \approx t' \text{\quad yen lan mung yen \quad} \eq[t][t'] \in \Gamma^*
  \]
\end{defn}

\begin{prop}
\ollabel{prop:approx-equiv}
Relasi $\approx$ nduweni sipat-sipat ing ngisor iki:
\begin{enumerate}
\item $\approx$ iku refleksif.
\item $\approx$ iku simetris.
\item $\approx$ iku transitif.
\item Yen $t \approx t'$, $f$ minangka !!a{function}, lan $t_1$,
  \dots, $t_{i-1}$, $t_{i+1}$, \dots, $t_n$ minangka term katutup,
  mula
\[
\Atom{f}{t_1,\dots, t_{i-1}, t, t_{i+1}, \dots, t_n} \approx
\Atom{f}{t_1,\dots, t_{i-1}, t', t_{i+1}, \dots, t_n}.
\]
\item Yen $t \approx t'$, $R$ minangka !!a{predicate}, lan $t_1$,
  \dots, $t_{i-1}$, $t_{i+1}$, \dots, $t_n$ minangka term katutup,
  mula
\begin{multline*}
\Atom{R}{t_1,\dots, t_{i-1}, t, t_{i+1}, \dots, t_n} \in \Gamma^* \text{ yen lan mung yen } \\
\Atom{R}{t_1,\dots, t_{i-1}, t', t_{i+1}, \dots, t_n} \in \Gamma^*.
\end{multline*}
\end{enumerate}
\end{prop}

\begin{proof}
Amarga $\Gamma^*$ konsisten lan !!{complete}, $\eq[t][t'] \in
\Gamma^*$ yen lan mung yen $\Gamma^* \Proves \eq[t][t']$. Mula cukup
nuduhake pratelan-pratelan ing ngisor iki:
\begin{enumerate}
\item $\Gamma^* \Proves \eq[t][t]$ kanggo kabeh term katutup~$t$.
\item Yen $\Gamma^* \Proves \eq[t][t']$, mula $\Gamma^* \Proves
  \eq[t'][t]$.
\item Yen $\Gamma^* \Proves \eq[t][t']$ lan $\Gamma^* \Proves
  \eq[t'][t'']$, mula $\Gamma^* \Proves \eq[t][t'']$.
\item Yen $\Gamma^* \Proves \eq[t][t']$, mula
\[
\Gamma^* \Proves
\eq[\Atom{f}{t_1,\dots,t_{i-1},t,t_{i+1},\dots,t_n}][\Atom{f}{t_1,\dots,t_{i-1},t',t_{i+1},\dots,t_n}]
\]
kanggo saben !!{function} $n$-panggonan~$f$ lan term katutup
$t_1$, \dots, $t_{i-1}$, $t_{i+1}$, \dots,~$t_n$. Cathetan koreksi
sumber OLPL-050: koma rangkep sawise term kanthi indeks i tambah siji ing
rumus sumber dibusak.
\item Yen $\Gamma^* \Proves \eq[t][t']$ lan
$\Gamma^* \Proves
\Atom{R}{t_1,\dots,t_{i-1},t,t_{i+1},\dots,t_n}$, mula
$\Gamma^* \Proves \Atom{R}{t_1,\dots,t_{i-1},t',t_{i+1},\dots,t_n}$
kanggo saben !!{predicate} $n$-panggonan~$R$ lan term katutup
$t_1$, \dots, $t_{i-1}$, $t_{i+1}$, \dots,~$t_n$.
\end{enumerate}
\end{proof}

\begin{prob}
Rampungna bukti~\olref[fol][com][ide]{prop:approx-equiv}.
\end{prob}

\begin{defn}
Upamana $\Gamma^*$ minangka himpunan konsisten lan !!{complete} ing
basa~$\Lang L$, $t$ minangka term katutup, lan $\approx$ kaya ing
definisi sadurunge. Banjur:
\[
\equivrep{t}{\approx} = \Setabs{t'}{t'\in \Trm[L], t \approx t'}
\]
lan $\equivclass{\Trm[L]}{\approx} = \Setabs{\equivrep{t}{\approx}}{t \in \Trm[L]}$.
\end{defn}

\begin{defn}
\ollabel{defn:term-model-factor}
Upamana $\Struct M = \Struct M(\Gamma^*)$ minangka model term
kanggo~$\Gamma^*$ saka \olref[mod]{defn:termmodel}. Banjur
$\Struct{\equivclass{M}{\approx}}$ yaiku !!{structure} ing ngisor iki:
\begin{enumerate}
\item $\Domain{\equivclass{M}{\approx}} = \equivclass{\Trm[L]}{\approx}$.
\item $\Assign{c}{\equivclass{M}{\approx}} = \equivrep{c}{\approx}$
\item $\Assign{f}{\equivclass{M}{\approx}}(\equivrep{t_1}{\approx}, \dots,
  \equivrep{t_n}{\approx}) = \equivrep{\Atom{f}{t_1,\dots, t_n}}{\approx}$
\item $\tuple{\equivrep{t_1}{\approx}, \dots, \equivrep{t_n}{\approx}} \in
  \Assign{R}{\equivclass{M}{\approx}}$ yen lan mung yen
  $\Sat{M}{\Atom{R}{t_1,\dots, t_n}}$, yaiku yen lan mung yen
  $\Atom{R}{t_1,\dots, t_n} \in \Gamma^*$.
\end{enumerate}
\end{defn}

\begin{explain}
Elinga yen kita wis nemtokake $\Assign{f}{\equivclass{M}{\approx}}$
lan $\Assign{R}{\equivclass{M}{\approx}}$ kanggo unsur-unsur saka
$\equivclass{\Trm[L]}{\approx}$ kanthi ngrujuk unsur kasebut minangka
$\equivrep{t}{\approx}$, yaiku lumantar
\emph{wakil}~$t \in \equivrep{t}{\approx}$. Kita kudu mesthekake yen
definisi iki ora gumantung marang pilihan wakil kasebut, yaiku yen
pilihan liyane~$t'$ nemtokake kelas ekuivalensi kang padha
($\equivrep{t}{\approx} = \equivrep{t'}{\approx}$), definisine menehi
asil kang padha. Umpamane, yen $R$ minangka !!{predicate} kanthi siji
panggonan, klausa pungkasan ing definisi nyatakake yen
$\equivrep{t}{\approx} \in \Assign{R}{\equivclass{M}{\approx}}$ yen lan
mung yen $\Sat{M}{\Atom{R}{t}}$. Yen kanggo term liyane~$t'$ kanthi
$t \approx t'$ nanging $\Sat/{M}{\Atom{R}{t'}}$, definisi kasebut
bakal mbutuhake $\equivrep{t'}{\approx} \notin
\Assign{R}{\equivclass{M}{\approx}}$. Cathetan koreksi sumber OLPL-052:
makro Sat kanthi garis miring wis nyatakake yen predikat luput; conto
kasebut mung mbutuhake wakil kapindho mawa tandha prima, dudu wakil
kapisan. Yen $t \approx t'$, mula
$\equivrep{t}{\approx} = \equivrep{t'}{\approx}$, nanging ora mungkin
$\equivrep{t}{\approx} \in \Assign{R}{\equivclass{M}{\approx}}$ lan
$\equivrep{t}{\approx} \notin
\Assign{R}{\equivclass{M}{\approx}}$ kalorone. Nanging
\olref{prop:approx-equiv} njamin yen iki ora bisa kedadeyan.
\end{explain}

\begin{prop}
$\Struct{\equivclass{M}{\approx}}$ ditemtokake kanthi becik, yaiku yen
$t_1$, \dots, $t_n$, $t_1'$, \dots, $t_n'$ minangka term katutup, lan
$t_i \approx t_i'$, mula
\begin{enumerate}
\item $\equivrep{\Atom{f}{t_1,\dots, t_n}}{\approx} =
    \equivrep{\Atom{f}{t_1',\dots, t_n'}}{\approx}$, yaiku
  \[
  \Atom{f}{t_1,\dots, t_n} \approx \Atom{f}{t_1',\dots, t_n'}
  \]
  lan
\item $\Sat{M}{\Atom{R}{t_1,\dots, t_n}}$ yen lan mung yen
  $\Sat{M}{\Atom{R}{t_1',\dots, t_n'}}$, yaiku
  \[
    \Atom{R}{t_1,\dots, t_n} \in \Gamma^* \text{ yen lan mung yen }
    \Atom{R}{t_1',\dots, t_n'} \in \Gamma^*.
  \]
\end{enumerate}
\end{prop}

\begin{proof}
Iki ndherek saka \olref{prop:approx-equiv} kanthi induksi marang~$n$.
\end{proof}

Kaya ing perkara model term, sadurunge mbuktekake lemma bebener kita
mbutuhake lemma ing ngisor iki.

\begin{lem}
\ollabel{lem:val-in-termmodel-factored} Upamana $\Struct M = \Struct M
 (\Gamma^*)$, banjur kanggo saben term katutup
$\Value{t}{\equivclass{M}{\approx}} = \equivrep{t}
 {\approx}$. Cathetan koreksi sumber OLPL-051: domain lan bukti kang
dirujuk mbutuhake watesan term katutup.
\end{lem}
\begin{proof}
Buktine padha karo bukti \olref[mod]{lem:val-in-termmodel}.
\end{proof}

\begin{prob}
Rampungna bukti~\olref[fol][com][ide]{lem:val-in-termmodel-factored}.
\end{prob}

\begin{lem}
\ollabel{lem:truth}
$\Sat{\equivclass{M}{\approx}}{!A}$ yen lan mung yen $!A \in
\Gamma^*$ kanggo kabeh ukara~$!A$.
\end{lem}

\begin{proof}
Kanthi induksi marang~$!A$, kaya ing bukti \olref[mod]{lem:truth}.
Siji-sijine perkara kang mbutuhake kawigaten tambahan yaiku nalika
$!A \ident \eq[t][t']$.
\begin{align*}
\Sat{\equivclass{M}{\approx}}{\eq[t][t']} & \text{ yen lan mung yen } \equivrep{t}{\approx} = \equivrep{t'}{\approx}
\text{ (miturut definisi $\Struct{\equivclass{M}{\approx}}$)}\\
& \text{ yen lan mung yen } t \approx t' \text{ (miturut definisi $\equivrep{t}{\approx}$)}\\
& \text{ yen lan mung yen } \eq[t][t'] \in \Gamma^* \text{ (miturut definisi $\approx$).}
\end{align*}
\end{proof}

\begin{digress}
Elinga yen sanajan $\Struct{M(\Gamma^*)}$ mesthi !!{enumerable} lan
tanpa wates, $\Struct{\equivclass{M}{\approx}}$ bisa winates, amarga
bisa wae mung ana winates akeh kelas $\equivrep{t}{\approx}$. Iki
pancen samesthine, amarga $\Gamma$ bisa ngemot !!{sentence}s kang
mbutuhake saben !!{structure} panggonane ukara kasebut bener dadi
winates. Umpamane, $\lforall[x][\lforall[y][x = y]]$ minangka
!!{sentence} konsisten, nanging mung dicukupi ing !!{structure}s kanthi
!!a{domain} kang ngemot persis siji !!{element}.
\end{digress}

\end{document}
